\section{From the separating node to an equivariant flag eigensheaf}
\label{sec:nodal}

The parameter $\tau$ of Proposition~\ref{prop:nodal-parameter} lives on
the smooth projective generic curve $C/\Omega$ of the twisted
degeneration.  Proposition~\ref{prop:equivariant-unramified} therefore
supplies an equivariant unramified eigencomplex.  We now specialize
this complex to the node, retain the first normalization component,
and remove the enhancement of its flag.  The two issues are
nonvanishing on the prescribed node type and preservation of the
semilinear action when passing to a perverse eigensheaf.

\begin{proposition}\label{prop:charzero-equivariant}
Let $(C_1,x_1)/D$, $\sigma_1$, and $\varphi$ be the pointed curve,
dense lifted parameter, and semilinear automorphism of
Proposition~\ref{prop:lifted-parameter}.  There exist an integer
$m\geq 1$ and a nonzero locally constructible perverse geometric
$\E$-adic sheaf
\[
 M_{\Dbar}\quad\text{on}\quad
 \mathcal A_1=\Bun_{G,B,x_1}((C_1)_{\Dbar})
\]
with a $\varphi^m$-structure and the full coherent Hecke eigenstructure
for $\sigma_1$.  Every eigenisomorphism is compatible with this
structure and with the given $\varphi^m$-structure on $\sigma_1$.
Moreover, $\SS(M_{\Dbar})\subset\Lambda_{\mathrm{par}}$.
\end{proposition}

\subsection{The prescribed node type and its two flags}

Retain the fixed index $b_*$ and strictly dominant cocharacter
$\lambda$ from Section~\ref{sec:parameters}; thus
\[
 0<\langle\alpha,\lambda\rangle<b_*
 \qquad\text{for every positive root }\alpha.
\]
Put $S=S_{b_*}=\Spec\Dbar[[s_{b_*}]]$.  In the relative bundle stack
of $\mathcal C(b_*)/S$, remove from the closed fiber every inertia
type except the conjugacy class of $\lambda|_{\mu_{b_*}}$.  Denote
the resulting open stack by $\mathcal B$.  The condition is open
because the finitely many conjugacy classes of homomorphisms
$\mu_{b_*}\to G$ are locally constant on the residual gerbe.
The centralizer of the chosen type is $T$: the inequalities ensure
that no root is trivial on $\lambda(\mu_{b_*})$.
The Hom-stack theorem gives algebraicity and local finite
presentation \cite[Theorem~1.2]{HallRydh}.  Smoothness follows from
the vanishing of $H^2$ of the adjoint bundle: the curve is tame,
coarse pushforward is exact, and its coarse space has dimension one.
In particular $\mathcal B$ has smooth finite-type charts over $S$.
Its geometric generic fiber is $\Bun_G(C)$.

Let $B^-$ be the Borel opposite to $B$ with the same maximal torus
$T$, and define
\[
 \mathcal A_1^+=\Bun_{G,R_u(B),x_1}((C_1)_{\Dbar}),
 \qquad
 \mathcal A_2^+=\Bun_{G,R_u(B^-),x_2}((C_2)_{\Dbar}).
\]
An object of either stack includes an enhanced flag, that is, a
reduction to the indicated unipotent radical.  The calculation of
\cite[Lemma~8.2]{CT} gives
\begin{equation}\label{eq:nodal-type-quotient}
 \mathcal B_s\simeq
       [\mathcal A_1^+\times\mathcal A_2^+/T],
\end{equation}
where $T$ changes the two frames on the node gerbe simultaneously.
We recall the calculation to verify its compatibility with our
semilinear action and with the two different component curves.

On the first branch of the balanced chart, set $a=u^{b_*}$.
In a frame in which inertia is $\lambda$, a change of frame satisfies
\[
 h(\zeta u)=\lambda(\zeta)h(u)\lambda(\zeta)^{-1}.
\]
Conjugating to the invariant punctured-disc frame gives
$g(a)=\lambda(u)^{-1}h(u)\lambda(u)$.  If
$j=\langle\alpha,\lambda\rangle$ for a positive root, the positive
and negative root coordinates respectively transform as
\begin{equation}\label{eq:nodal-root-coordinate}
 u^j f(u^{b_*})\longmapsto f(a),
 \qquad
 u^{b_*-j}g(u^{b_*})\longmapsto a g(a).
\end{equation}
The toral coordinates are series in $a$.  These formulas identify
the full frame-change group with
$\{g(a)\in G(\Dbar[[a]]):g(0)\in B\}$; the assertion over arbitrary
test schemes follows by the same root-group factorization in the
formal big cell.  Evaluation $h\mapsto h(0)$ becomes the torus
projection of this Iwahori group.  Fixing the gerbe frame therefore
gives an enhanced $B$-flag.

On the second branch the node generator acts by
$v\mapsto\zeta^{-1}v$.  Its positively oriented coordinate generator
thus sees $-\lambda$, so the identical calculation gives $B^-$.
On both branches the fixed gerbe torsor has automorphism group $T$,
expressed using the common node generator.  Gluing the normalization bundles, with an
identification on their gerbes, now gives
\eqref{eq:nodal-type-quotient}.  The computation is branchwise and
requires no isomorphism between $C_1$ and $C_2$.

All these prescriptions use the split group, the group-scheme
cocharacter $\lambda$, and the balanced coordinates chosen over
$D$.  They consequently commute with $\varphi$, including its
action on $\mu_{b_*}$; no individual root of unity has been declared
fixed.  Thus \eqref{eq:nodal-type-quotient} is an equivariant
identification.  A modification away from the node preserves its
gerbe frames and acts on the corresponding factor.  This also
identifies every successive and multi-leg Hecke correspondence,
after pullback to $\mathcal A_1^+\times\mathcal A_2^+$, with its
componentwise counterpart.

\subsection{A criterion for nonzero nearby cycles}

Nearby cycles can vanish for a nonzero generic complex.  The input
that excludes this possibility here is the specialization-detection
theorem of \cite[Theorem~4.2]{CT}, with the following precise scope.

\begin{lemma}[Specialization detection]\label{lem:specialization-detection}
Let $S=\Spec A$ be an excellent complete strictly henselian trait
with algebraically closed residue field and $\ell$ invertible.
Let $\mathcal Y\to S$ be smooth and locally of finite type.  Its
special fiber is to be an unlevelled, ordinary Borel-level, or
enhanced Borel-level bundle stack on a smooth projective curve, or,
in characteristic zero, the simultaneous-torus quotient of two
enhanced-level stacks.  If the residue characteristic is positive,
assume that the residue field is an algebraic closure of a finite
field and impose the Lie-theoretic hypotheses of
Lemma~\ref{lem:lie-hypotheses}.

Assume that a smooth scheme $L\to S$ of relative dimension one
provides spherical Hecke legs and that $L_s$ contains a nonempty
open of each special-fiber curve on which modifications are allowed.
The bounded output maps to $\mathcal Y\times_S L$ must be
representable and proper; exact-relative-position strata must be
smooth over both input-and-leg and output-and-leg; their special
fibers must have the transverse-modification geometry of
\cite[Proposition~3.3]{CT}.  The Satake kernel on its open stratum is
constant with the relative Satake shift and twist.

Let $F$ be a locally bounded constructible complex on
$\mathcal Y_{\bar\eta}$ with lisse spherical Hecke eigenvalues on
$L_{\bar\eta}$.  Suppose each eigenvalue has a lisse lattice whose
finite reductions extend lisse after finite trait extensions,
compatibly with the special-fiber value, as in
Proposition~\ref{prop:equivariant-specialization}.
If, on some smooth finite-type chart $Z\to\mathcal Y$, the closure
of the nonzero-stalk locus of $F|_{Z_{\bar\eta}}$ meets $Z_s$, then
$\Psi F\ne0$.
\end{lemma}

Here and below $\Psi$ is geometric nearby cycles, without inertia
invariants.  The closure in the lemma may be computed after
descending its finitely many geometric generic components to a
finite extension of the trait.  Neither the lemma nor this
interpretation requires the complex $F$ itself to descend.
These are part of the statement and conventions of
\cite[Theorem~4.2]{CT}.

We will produce the required meeting by extending a bundle carrying
a nonzero generic stalk.  For $G=\SL_n$, the uniformization needed
for this step has an elementary description.  If $Y$ is a smooth
projective connected curve over an algebraically closed field and
$y\in Y$, the complement $Y\setminus\{y\}$ is a smooth affine curve.
Its coordinate ring is a Dedekind domain.  A rank-$n$ projective
module over that ring is isomorphic to the sum of a free module of
rank $n-1$ and its determinant.  Consequently every $\SL_n$-bundle
is trivial on this complement.  A free basis can be adjusted by a
unit so that its determinant agrees with the given trivialization.
This is the $\SL_n$ case of the affine-curve triviality used in
Drinfeld--Simpson uniformization and in
\cite[Theorem~1.1]{BelkaleFakhruddin}; see also
\cite{DrinfeldSimpson}.

In particular, any two $\SL_n$-bundles on $Y$ differ by a
modification at $y$.  Each such modification lies in some finite
Schubert bound.  In a family, the bounded Hecke space with fixed
reference input and fixed smooth section $y$ is proper over the
base trait.  A geometric generic modification therefore extends
after a finite trait extension: its point first descends to a
finite generic extension, and the valuative criterion then applies
over the corresponding DVR normalization of the complete trait.
The output remains unchanged near every
marking or node disjoint from $y$.

\subsection{Specializing on the chosen component}

Apply Proposition~\ref{prop:equivariant-unramified} to $(C,\tau)$.
The generic curve is smooth, projective, and connected of genus at
least two; its polarization comes from the projective smoothing.
Proposition~\ref{prop:nodal-parameter} supplies the continuous dense
parameter over a finite coefficient field and its equivariance.
We obtain a nonzero locally bounded constructible eigencomplex
$D_\tau$ on $\mathcal B_{\bar\eta}=\Bun_G(C)$, with its
$\varphi$-structure and coherent equivariant eigenmaps.

The trait $S_{b_*}$ is excellent, complete, and strictly henselian,
with algebraically closed characteristic-zero residue field; in
particular $\ell$ is invertible.  Use $L=\mathcal V$, the smooth
scheme locus of $\mathcal C(b_*)$, a relative curve with
$L_{\bar\eta}=C$ and $L_s=U_{1,\Dbar}\sqcup U_{2,\Dbar}$.
At these legs the bounded correspondences have the split affine
Grassmannian models in smooth frame charts.  Their output maps are
representable and proper, and their exact-position maps are smooth
over both bundle-and-leg projections.  These local models have the
normalized constant kernels on their open cells.  Modifications
there preserve the node type.  On the special fiber
\eqref{eq:nodal-type-quotient}, the cotangent description, cone
dimension bound, and transverse modifications are precisely the
torus-quotient case of \cite[Section~3 and Section~8.2]{CT}.
The required Lie-theoretic properties hold in characteristic zero.
Thus the geometric hypotheses of
Proposition~\ref{prop:equivariant-specialization} and
Lemma~\ref{lem:specialization-detection} hold for this family.

For their eigenvalue hypothesis, use the equivariant finite-torsor
tower of Proposition~\ref{prop:nodal-parameter}.  For each
representation, an invariant lattice has reductions extending
lisse over the smooth scheme locus after the extensions
$S_{b_r}\to S_{b_*}$.  Their specializations are the component
parameter systems, and all tensor maps and semilinear actions are
carried by the same tower.  Hence
\begin{equation}\label{eq:nodal-nearby}
 F_s:=\Psi D_\tau
\end{equation}
is locally bounded constructible and carries the coherent component
eigenstructure and its $\varphi$-action.

To prove $F_s\ne0$ on this particular type, first choose a point of
$\mathcal A_1^+\times\mathcal A_2^+$.  A smooth chart of
$\mathcal B$ through its image, followed by henselian lifting,
gives a reference bundle $P_0$ over $S$ with the prescribed node
type.  Lift also a point of $L_s$ to a section $y$ of $L/S$.
Choose a bundle $P$ on $C/\Omega$ where $D_\tau$ has nonzero stalk.
Such a point exists on a finite-type chart because $D_\tau$ is
nonzero and locally bounded constructible.

The Dedekind-domain argument above identifies $P$ and $(P_0)_\Omega$
off $y_\Omega$.  Extend the resulting bounded modification after a
finite trait extension.  Its output lies in $\mathcal B$, since its
node restriction still equals that of $P_0$.  Take a smooth
finite-type chart $Z\to\mathcal B$ through its special value and a
residue-field point of the chart above it.  If $S'$ is the extended
trait, the morphism $Z\times_{\mathcal B}S'\to S'$ is smooth;
henselian lifting of that point gives a section, after a further
finite extension if necessary.  Its geometric generic
point has the chosen nonzero stalk.  This section witnesses a
meeting of the closure in Lemma~\ref{lem:specialization-detection}
with the special fiber.  Thus
\begin{equation}\label{eq:nodal-nonzero}
 F_s\ne0\quad\text{on }[\mathcal A_1^+\times\mathcal A_2^+/T].
\end{equation}
Finite extensions in this argument only locate the closure inside
the fixed geometric generic fiber; they impose no new descent
condition on $D_\tau$ or its semilinear structure.  They need not
carry $\varphi$: geometric nearby cycles are unchanged by this
finite-trait replacement, and the equivariant object
\eqref{eq:nodal-nearby} was already constructed over $S$.

Ordinary pullback of $F_s$ to
$\mathcal A_1^+\times\mathcal A_2^+$ is nonzero, because this map is
smooth and surjective.  Choose a $\Dbar$-point $(a_1,a_2)$ with
nonzero stalk.  The bundle and enhanced flag defining $a_2$ descend
to a finite extension $D'/D$: they are of finite presentation and
$\Dbar/D$ is algebraic.  Pass to a finite normal extension
containing $D'$.  Since $\varphi$ fixes $D$, its action on this
normal extension has finite order.  Some power $\varphi^m$ fixes
it pointwise, and the descended model then supplies an
identification $\varphi^{m*}a_2\simeq a_2$.  We henceforth use this
power as generator.

Restricting along $\mathcal A_1^+\times\{a_2\}$ gives a nonzero
locally bounded constructible complex $F_1^+$ on $\mathcal A_1^+$,
with a $\varphi^m$-structure.  The quotient description of the Hecke
correspondences and proper base change for each bound give
\begin{equation}\label{eq:nodal-enhanced-eigen}
 \Hecke_{I,(V_i)}(F_1^+)
   \simeq F_1^+\boxtimes
                 \mathop{\boxtimes}_{i\in I}(V_i)_{\sigma_1}
\end{equation}
for every finite leg set $I$.  The same comparisons apply on the
successive-modification spaces and on every fusion diagonal.
They are natural under semilinear transport, so all these
eigenmaps and their coherence constraints are equivariant.
Ordinary restriction here is sufficient; no perversity assertion
is made for $F_1^+$.

\subsection{Perverse cohomology and removal of the enhancement}

To pass from $F_1^+$ to an ordinary-level perverse eigensheaf, we
will take perverse cohomology, push forward along the torsor that
forgets the enhancement, and take perverse cohomology once more.
The first truncation supplies nilpotent singular support, which
controls monodromy along the enhancement torus.  Unipotence of that
monodromy will ensure that the direct image is nonzero; the second
truncation then gives the required perverse object.

The geometric inputs for these steps are
\cite[Propositions~6.1 and 7.5]{CT}.  Their scope is a smooth
projective complex curve with one marked
point and a split simple simply connected group.  Proposition~6.1
applies at either ordinary or enhanced Borel level: for a locally
bounded constructible geometric adic complex with a coherent
eigenstructure for a Zariski-dense parameter, every perverse
cohomology inherits the full eigenstructure and nilpotent singular
support, and some perverse cohomology is nonzero if the complex is.
Proposition~7.5 removes the enhancement when the boundary monodromy
is unipotent.  Our group $\SL_n$ meets the group assumptions, and
Proposition~\ref{prop:lifted-parameter} proves both density and the
required unipotent boundary monodromy for $\sigma_1$.

Choose an abstract field isomorphism $\Dbar\simeq\C$ to apply these
geometric statements.  Such an isomorphism exists: both fields
are algebraically closed of characteristic zero and have
transcendence degree $2^{\aleph_0}$ over $\Q$.  We will use their
Betti arguments to establish geometric properties, while making
all sheaf operations and induced maps in the geometric adic
category over $\Dbar$.

First choose $j$ such that $Q={}^pH^j(F_1^+)\ne0$.
Here is why its eigenstructure retains the semilinear action.
For a finite set $I$, put $d=|I|$.  The proof of
\cite[Proposition~6.1, Section~6.3]{CT} constructs, on the objects
under consideration, the functorial adic comparison
\begin{equation}\label{eq:nodal-truncation}
 {}^pH^j\bigl(\Hecke_{I,(V_i)}(F_1^+)[d]\bigr)
 \simeq \Hecke_{I,(V_i)}({}^pH^jF_1^+)[d].
\end{equation}
The Betti calculation proves the needed perverse bounds; the
comparison itself follows from the universal property of the
adic truncation triangles.  On the eigenvalue side, put
$\mathcal L_I=\boxtimes_{i\in I}(V_i)_{\sigma_1}$.  Since this
sheaf is lisse on the smooth $d$-dimensional leg space, exterior
product with $\mathcal L_I[d]$ is perverse exact, so
\[
 {}^pH^j(F_1^+\boxtimes\mathcal L_I[d])
       \simeq Q\boxtimes\mathcal L_I[d].
\]
Applying these comparisons to \eqref{eq:nodal-enhanced-eigen}
and canceling the common shift $[d]$ gives the relative-normalized
eigenmaps for $Q$.

For a collision $\Delta:U_1^J\to U_1^I$, the corresponding
truncation comparison uses
$(\id\times\Delta)^*[|J|-|I|]$ on these locally constant leg
families.  The comparisons for successive Hecke operations use the
total dimension of their retained leg space, counting a shared
leg only once.  Thus \eqref{eq:nodal-truncation} carries the entire
unit, convolution, permutation, and fusion diagrams.  Pullback by
$\varphi^m$ preserves the perverse structure and its truncation
maps, so these adic comparisons commute with it.  No compatibility
between a chosen complex realization and $\varphi$ is needed.

Now let
\[
 q:\mathcal A_1^+\longrightarrow\mathcal A_1
\]
forget the enhancement.  This is a $T$-torsor of relative dimension
$r=n-1$.  Set $F=q_!Q$.  We need the intermediate assertion in the
proof of \cite[Proposition~7.5]{CT}: this particular direct image is
nonzero.  We recall its mechanism, since a general torus direct
image does not preserve nonvanishing.

At enhanced level, a cotangent residue lies in $\operatorname{Lie}B$.
In an adapted local frame write a generically nilpotent Higgs field
as $\theta=b(t)\,dt/t$, with $b(t)$ regular.  The positive-degree
invariant polynomials vanish on $b(t)$ and hence on the residue
$b(0)$; the toral projection of this residue is therefore zero.
Thus the enhanced nilpotent cone is the smooth cotangent pullback
of the ordinary-level cone and annihilates tangent directions to
the $T$-fibers \cite[Section~3.1]{CT}.  The microlocal criterion of
\cite[Lemma~7.1]{CT} now makes the Betti realization of $Q$
monodromic along this torsor: locally on the base it is constructible
for a product stratification with the whole torus as second factor.
In particular its cohomology sheaves on each fiber are local systems.

For a representation $V$, let $\chi_V$ be its character on the
dual torus $\widehat T$.  The central-sheaf calculation, based on
Gaitsgory's construction \cite{GaitsgoryCentral} and proved in the
needed form in \cite[Lemma~7.4 and proof of Proposition~7.5]{CT},
gives, for
every $V\in\Rep(\Gd)$,
\begin{equation}\label{eq:nodal-character}
 \chi_V(\text{enhancement monodromy})
       =\operatorname{Tr}(\text{boundary monodromy on }V_{\sigma_1})
          \,\id_Q
       = (\dim V)\,\id_Q.
\end{equation}
The last equality uses unipotence at $x_1$.  On any stalk
cohomology of a torsor fiber, the commuting monodromies have joint
generalized eigencharacters, viewed as points of the dual torus
$\widehat T(\E)$.  Character functions of representations span the
Weyl-invariant regular functions on $\widehat T$ and separate its
Weyl orbits.  Equation
\eqref{eq:nodal-character} therefore puts each eigencharacter in
the orbit of $1$, which consists of $1$ alone.  All enhancement
monodromies on these finite-dimensional spaces are unipotent.

Choose a fiber $T=q^{-1}(a)$ where $Q_T:=Q|_T$ is nonzero, and let
$j_0$ be the largest index with $\mathcal H^{j_0}(Q_T)\ne0$.
Its fiber $W$ is a nonzero finite-dimensional representation of
$\Gamma=X_*(T)$ with commuting unipotent operators.  The
augmentation ideal of $\E[\Gamma]$ acts nilpotently on this space,
so the coinvariants $W_\Gamma$ are nonzero.  Poincar\'e duality gives
\[
 H_c^{2r}(T,\mathcal H^{j_0}(Q_T))
       \simeq W_\Gamma(-r)\ne0.
\]
In the compact-support hypercohomology spectral sequence this term
has no incoming differential, by maximality of $j_0$, and no
outgoing differential, because $H_c^i(T,-)=0$ for $i>2r$.
Consequently $H_c^{2r+j_0}(T,Q_T)\ne0$.  Comparison and base change
for the adic functor $q_!$ imply $F_a\ne0$.  This argument uses a
nilpotence bound only on the chosen finite-dimensional space.

The morphism $q$ has finite type and fixed relative dimension, so
$F$ is locally bounded constructible.  Away from $x_1$, Hecke
modifications transport an enhanced flag as well as its ordinary
flag.  Proper base change on a bounded Hecke correspondence and
the projection formula therefore give the natural comparison
\begin{equation}\label{eq:nodal-torus-hecke}
 \Hecke_{I,(V_i)}(q_!Q)
 \simeq (q\times\id_{U_1^I})_!\Hecke_{I,(V_i)}(Q)
 \simeq q_!Q\boxtimes
                  \mathop{\boxtimes}_{i\in I}(V_i)_{\sigma_1}.
\end{equation}
These comparisons also hold on the successive and collision
correspondences.  Hence they preserve every coherence constraint.
Since $q$ is defined over $D$, its direct image and these natural
comparison maps carry the $\varphi^m$-structure.

Finally choose a nonzero perverse cohomology
$M_{\Dbar}={}^pH^a(F)$.  Proposition~6.1 of \cite{CT} applies now
at ordinary Borel level, and the same adic truncation argument
\eqref{eq:nodal-truncation} proves equivariance of all the induced
eigenmaps.  It also gives nilpotent singular support.  At ordinary
flag level this is $\Lambda_{\mathrm{par}}$.
This completes the proof of Proposition~\ref{prop:charzero-equivariant}.
